\phantomsection
\section*{LỜI MỞ ĐẦU}
\addcontentsline{toc}{section}{LỜI MỞ ĐẦU}
Giao hàng đúng hẹn là một yêu cầu quan trọng trong chuỗi cung ứng. Việc
nhận diện sớm đơn có nguy cơ giao trễ có thể hỗ trợ theo dõi và phân bổ
nguồn lực, nhưng cảnh báo sai hoặc bỏ sót đều cần được xem xét. Từ bối
cảnh đó, nhóm lựa chọn đề tài dự đoán nguy cơ giao trễ để vận dụng kiến
thức học phần vào một bài toán phân loại dữ liệu bảng.

Báo cáo sử dụng DataCo Smart Supply Chain, với K-Nearest Neighbors và
Decision Tree là hai thuật toán chính. Quy trình gồm kiểm toán dữ liệu,
lựa chọn đặc trưng, tiền xử lý, chia nhóm theo mã đơn, kiểm định chéo và
đánh giá. Hai luồng A/B cùng các baseline được đối chiếu bằng kết quả
đã lưu trong repository. Các kết luận phân biệt rõ dòng hàng với đơn
hàng và giới hạn của tập con thực nghiệm, không suy rộng thành hiệu quả
vận hành đã được chứng minh.

Báo cáo gồm sáu chương theo thứ tự: đặt vấn đề và phân tích yêu cầu;
thiết kế giải pháp; phân tích mã nguồn; thử nghiệm và kết quả; thảo luận
và so sánh; tổng kết và hướng phát triển.

\phantomsection
\section*{PHÂN CÔNG NHIỆM VỤ}
\addcontentsline{toc}{section}{PHÂN CÔNG NHIỆM VỤ}
\begin{table}[!htbp]\centering\small
\caption{Phân công nhiệm vụ dự kiến, chờ hai thành viên xác nhận.}\label{tab:team-assignment}
\begin{tabularx}{\linewidth}{p{3.6cm}Y}\toprule
\rowcolor{reportheader} \theadcell{Họ tên thành viên} & \theadcell{Nhiệm vụ được phân công}\\\midrule
Vũ Quốc Đoàn\newline 1871020143 & Nghiên cứu bài toán, dataset; phân tích EDA; tiền xử lý và lựa chọn đặc trưng; phân tích luồng A/B; tham gia hoàn thiện báo cáo.\\[7pt]
Trần Đại Hiệp\newline 1871020223 & Nghiên cứu K-NN và Decision Tree; xây dựng pipeline huấn luyện; cross-validation và đánh giá; phân tích kết quả, ma trận nhầm lẫn và lỗi; tham gia hoàn thiện báo cáo.\\\bottomrule
\end{tabularx}
\end{table}
